\documentclass[10pt,a4paper]{amsart}
\usepackage{amsmath,amssymb,amsthm,graphicx}
% Retained minimal preamble for the Stacks Project chapter more-morphisms.tex (More on Morphisms).
% Verbatim excerpt of the chapter preamble: Xy-pic setup (native preamble lines 14-18) and the native
% theorem declarations with the mathematical macros (native preamble lines 154-217). Copied mechanically;
% nothing added, removed or redefined. The wrapper class and the amsmath/amssymb/amsthm/graphicx base
% packages are supplied by the assembler, so the class/externaldocument part of the native preamble is omitted.
% For commutative diagrams we use Xy-pic
\usepackage[all]{xy}

% We use 2cell for 2-commutative diagrams.
\xyoption{2cell}
% Theorem environments.
%
\theoremstyle{plain}
\newtheorem{theorem}[subsection]{Theorem}
\newtheorem{proposition}[subsection]{Proposition}
\newtheorem{lemma}[subsection]{Lemma}

\theoremstyle{definition}
\newtheorem{definition}[subsection]{Definition}
\newtheorem{example}[subsection]{Example}
\newtheorem{exercise}[subsection]{Exercise}
\newtheorem{situation}[subsection]{Situation}

\theoremstyle{remark}
\newtheorem{remark}[subsection]{Remark}
\newtheorem{remarks}[subsection]{Remarks}

\numberwithin{equation}{subsection}

% Macros
%
\def\lim{\mathop{\mathrm{lim}}\nolimits}
\def\colim{\mathop{\mathrm{colim}}\nolimits}
\def\Spec{\mathop{\mathrm{Spec}}}
\def\Hom{\mathop{\mathrm{Hom}}\nolimits}
\def\Ext{\mathop{\mathrm{Ext}}\nolimits}
\def\SheafHom{\mathop{\mathcal{H}\!\mathit{om}}\nolimits}
\def\SheafExt{\mathop{\mathcal{E}\!\mathit{xt}}\nolimits}
\def\Sch{\mathit{Sch}}
\def\Mor{\mathop{\mathrm{Mor}}\nolimits}
\def\Ob{\mathop{\mathrm{Ob}}\nolimits}
\def\Sh{\mathop{\mathit{Sh}}\nolimits}
\def\NL{\mathop{N\!L}\nolimits}
\def\CH{\mathop{\mathrm{CH}}\nolimits}
\def\proetale{{pro\text{-}\acute{e}tale}}
\def\etale{{\acute{e}tale}}
\def\QCoh{\mathit{QCoh}}
\def\Ker{\mathop{\mathrm{Ker}}}
\def\Im{\mathop{\mathrm{Im}}}
\def\Coker{\mathop{\mathrm{Coker}}}
\def\Coim{\mathop{\mathrm{Coim}}}

% Boxtimes
%
\DeclareMathSymbol{\boxtimes}{\mathbin}{AMSa}{"02}

%
% Macros for moduli stacks/spaces
%
\def\QCohstack{\mathcal{QC}\!\mathit{oh}}
\def\Cohstack{\mathcal{C}\!\mathit{oh}}
\def\Spacesstack{\mathcal{S}\!\mathit{paces}}
\def\Quotfunctor{\mathrm{Quot}}
\def\Hilbfunctor{\mathrm{Hilb}}
\def\Curvesstack{\mathcal{C}\!\mathit{urves}}
\def\Polarizedstack{\mathcal{P}\!\mathit{olarized}}
\def\Complexesstack{\mathcal{C}\!\mathit{omplexes}}
% \Pic is the operator that assigns to X its picard group, usage \Pic(X)
% \Picardstack_{X/B} denotes the Picard stack of X over B
% \Picardfunctor_{X/B} denotes the Picard functor of X over B
\def\Pic{\mathop{\mathrm{Pic}}\nolimits}
\def\Picardstack{\mathcal{P}\!\mathit{ic}}
\def\Picardfunctor{\mathrm{Pic}}
\def\Deformationcategory{\mathcal{D}\!\mathit{ef}}
\usepackage{cleveref}
\crefname{originallemma}{original lemma}{original lemmas}
\Crefname{originallemma}{Original lemma}{Original lemmas}
\makeatletter
\newlabel{lemma-hom-from-finite-locally-free-into-affine}{{0BL3}{source}{}{sourcepointer-1}{}}
\newlabel{lemma-hom-from-finite-locally-free-into-affine@cref}{{[originallemma][0][]0BL3}{[1][0][]source}{}{}{}}
\makeatother
\crefname{originalreference}{original reference}{original references}
\Crefname{originalreference}{Original reference}{Original references}
\makeatletter
\newlabel{schemes-lemma-glue-functors}{{01JJ}{source}{}{sourcepointer-2}{}}
\newlabel{schemes-lemma-glue-functors@cref}{{[originalreference][0][]01JJ}{[1][0][]source}{}{}{}}
\makeatother
\title{The morphism functor from a finite locally free scheme is representable by a scheme}
\author{The Stacks Project Authors}
\date{stacks-project:more-morphisms.tex:local-snapshot-20261011; GNU FDL 1.2 or later}
\begin{document}
\maketitle
\noindent\textit{Editorial note.} Counted span is the author's own proof environment, source lines 20213-20254, immediately following the statement at 20200-20211.\par\medskip
\noindent\textit{Editorial note.} One same-chapter lemma and one cross-chapter glueing lemma are recorded in assembly.external\_references.\par\medskip
\noindent\textit{Editorial note.} The statement's condition (2) is used exactly once, at the covering step inside the span.\par\medskip
\noindent\textit{Allowed standard prerequisites (editorial).}\begin{itemize}
\item Representability when the source is finite free over an affine base. chapter lemma lemma-hom-from-finite-locally-free-into-affine, cited in the proof. Provides the local model: over an affine piece of S the functor is representable by an affine scheme.
\item Glueing functors. Schemes, lemma schemes-lemma-glue-functors, cited in the proof. Glues the local representing objects along the affine opens of S into a scheme representing the whole functor.
\item Affineness of finitely many points in a fibre. hypothesis (2) of the statement, used at proof lines 20206-20208 and 20241-20247. Ensures that any finite set of points in a fibre of X over S lies in an affine open, which is exactly what the local representability needs to cover the functor.
\end{itemize}\medskip
\ifdefined\hypertarget\hypertarget{sourcepointer-1}{}\fi
\noindent\textit{Source reference.} Original lemma 0BL3: sources/stacks-more-morphisms/pinned-original-source.tex, source lines 20171--20171. Navigation pointer only: the counted proof cites this statement by \textbackslash{}ref as a standing input; the deduction the counted proof has to make is carried out inside the counted span, and no part of that deduction is delegated to this pointer.\par\medskip
\ifdefined\hypertarget\hypertarget{sourcepointer-2}{}\fi
\noindent\textit{Source reference.} Original reference 01JJ: focused-seeds/geometry/S9-morph-hom-locally-free-representable/ref-tags.txt, source lines 4--4. Navigation pointer only: the counted proof cites this statement by \textbackslash{}ref as a standing input; the deduction the counted proof has to make is carried out inside the counted span, and no part of that deduction is delegated to this pointer.\par\medskip
\begin{lemma}
\label{lemma-hom-from-finite-locally-free-representable}
Let $Z \to S$ and $X \to S$ be morphisms of schemes.
Assume
\begin{enumerate}
\item $Z \to S$ is finite locally free, and
\item for all $(s, x_1, \ldots, x_d)$ where $s \in S$ and
$x_1, \ldots, x_d \in X_s$ there exists an affine open $U \subset X$
with $x_1, \ldots, x_d \in U$.
\end{enumerate}
Then $\mathit{Mor}_S(Z, X)$ is representable by a scheme.
\end{lemma}

\begin{proof}
Consider the set $I$ of pairs $(U, V)$ where $U \subset X$ and $V \subset S$
are affine open and $U \to S$ factors through $V$. For $i \in I$ denote
$(U_i, V_i)$ the corresponding pair. Set
$F_i = \mathit{Mor}_{V_i}(Z_{V_i}, U_i)$.
It is immediate that $F_i$ is a subfunctor of $\mathit{Mor}_S(Z, X)$.
Then we claim that conditions
(2)(a), (2)(b), (2)(c) of
Schemes, Lemma \ref{schemes-lemma-glue-functors} which proves the lemma.

\medskip\noindent
Condition (2)(a) follows from
Lemma \ref{lemma-hom-from-finite-locally-free-into-affine}.

\medskip\noindent
To check condition (2)(b) consider $T/S$ and $b \in \mathit{Mor}_S(Z, X)(T)$.
Thinking of $b$ as a morphism $T \times_S Z \to X$ we find an open
$b^{-1}(U_i) \subset T \times_S Z$. Clearly, $b \in F_i(T)$
if and only if $b^{-1}(U_i) = T \times_S Z$. Since the projection
$p : T \times_S Z \to T$ is finite hence closed, the set
$U_{i, b} \subset T$ of points $t \in T$ with
$p^{-1}(\{t\}) \subset b^{-1}(U_i)$ is open.
Then $f : T' \to T$ factors through $U_{i, b}$ if and only
if $b \circ f \in F_i(T')$ and we are done checking (2)(b).

\medskip\noindent
Finally, we check condition (2)(c) and this is where our condition
on $X \to S$ is used. Namely, consider
$T/S$ and $b \in \mathit{Mor}_S(Z, X)(T)$.
It suffices to prove that every $t \in T$
is contained in one of the opens $U_{i, b}$ defined
in the previous paragraph.
This is equivalent to the condition that
$b(p^{-1}(\{t\})) \subset U_i$ for some $i$
where $p : T \times_S Z \to T$ is the projection and
$b : T \times_S Z \to X$ is the given morphism.
Since $p$ is finite, the set $b(p^{-1}(\{t\})) \subset X$
is finite and contained in the fibre of $X \to S$ over
the image $s$ of $t$ in $S$.
Thus our condition on $X \to S$ exactly shows a
suitable pair exists.
\end{proof}

\end{document}
